%%----------------------------------------------------------------------------%%
%% NC State required accessible ETD setup
%%----------------------------------------------------------------------------%%
%% This file contains settings that should remain active for the final
%% accessible NC State thesis/dissertation build.
%%
%% Do not put draft-only design choices here. Those belong in
%% optional-accessible.tex.
%%----------------------------------------------------------------------------%%

%%----------------------------------------------------------------------------%%
%% Link colors and shared colors
%% The soft colors are defined here because optional-accessible.tex may use
%% them when the optional theorem-card style is enabled.
%%----------------------------------------------------------------------------%%

\usepackage{xcolor}

\definecolor{NCSULinkBlue}{HTML}{005389}
\definecolor{NCSUReynoldsRed}{HTML}{990000}

\definecolor{NCSUSoftBlue}{HTML}{EEF6FB}
\definecolor{NCSUSoftBlueTitle}{HTML}{D9ECF7}
\definecolor{NCSUBlueBorder}{HTML}{2F5D7C}

\definecolor{NCSUSoftGreen}{HTML}{F0F7F2}
\definecolor{NCSUSoftGreenTitle}{HTML}{DCEFE2}
\definecolor{NCSUGreenBorder}{HTML}{3F6F4F}

\definecolor{NCSUBoxGray}{HTML}{F6F6F6}
\definecolor{NCSUBoxGrayTitle}{HTML}{EDEDED}

%%----------------------------------------------------------------------------%%
%% Hyperlinks, PDF bookmarks, and PDF metadata
%%----------------------------------------------------------------------------%%
%% pdfdisplaydoctitle=true tells PDF viewers to display the document title
%% instead of the file name when possible.
%%
%% The visible title page is not enough for accessibility checkers. The PDF also
%% needs a title in the document properties. Do not edit the student name or
%% dissertation title here. Those values are set in student-information.tex.

\usepackage[
  unicode=true,
  bookmarksnumbered=true,
  bookmarksopen=true,
  linktoc=all,
  colorlinks=true,
  allcolors=black,
  pdfdisplaydoctitle=true,
  pdftitle={\NCSUThesisTitle},
  pdfauthor={\NCSUStudentFullName},
  pdfsubject={\NCSUPDFSubject},
  pdfkeywords={\NCSUPDFKeywords},
  pdfcreator={LuaLaTeX with LaTeX PDF management and tagpdf}
]{hyperref}

\usepackage{bookmark}

%%----------------------------------------------------------------------------%%
%% Microtypography
%%----------------------------------------------------------------------------%%
%% Microtype improves spacing and line breaks. This is mainly a readability
%% improvement and should not change the structure of the document.

\usepackage{microtype}

%%----------------------------------------------------------------------------%%
%% NC State Graduate School ETD formatting
%%----------------------------------------------------------------------------%%
%% NC State requires most text to be the same font size as the body text.
%% Section-level heading formatting is defined directly in ncsuthesis.cls so
%% that the template does not depend on the obsolete sectsty package.

\captionsetup{
  font=normalsize,
  labelfont=bf
}

%%----------------------------------------------------------------------------%%
%% Base theorem-like environments
%%----------------------------------------------------------------------------%%
%% The demonstration chapters and many mathematics dissertations use these
%% environments. The base definitions are intentionally plain and suitable
%% for the final ETD:
%%   - no tcolorbox,
%%   - no amsthm theorem-like tagging structure,
%%   - no colored backgrounds.
%%
%% Optional visual changes to these environments belong in optional-accessible.tex.

\newcounter{ncsuStatement}[chapter]
\renewcommand{\thencsuStatement}{\thechapter.\arabic{ncsuStatement}}

\newcommand{\NCSUStatementTitle}[3]{%
  \textbf{#1~#2%
  \if\relax\detokenize{#3}\relax
  \else
    : #3%
  \fi.}%
}

%% Arguments 3 and 4 are reserved for the optional visual style. They are
%% intentionally unused by the plain required style below.
\newcommand{\NCSUStatementBegin}[4]{%
  \refstepcounter{ncsuStatement}%
  \par\addvspace{12pt}%
  \noindent\NCSUStatementTitle{#1}{\thencsuStatement}{#2}\quad
  \begingroup
  \ignorespaces
}

\newcommand{\NCSUStatementEnd}{%
  \par
  \endgroup
  \par\addvspace{12pt}%
}

\DeclareDocumentEnvironment{definition}{O{}}{%
  \NCSUStatementBegin{Definition}{#1}{NCSUSoftGreenTitle}{NCSUGreenBorder}%
}{%
  \NCSUStatementEnd
}

\DeclareDocumentEnvironment{theorem}{O{}}{%
  \NCSUStatementBegin{Theorem}{#1}{NCSUSoftBlueTitle}{NCSUBlueBorder}%
  \itshape
}{%
  \NCSUStatementEnd
}

\DeclareDocumentEnvironment{lemma}{O{}}{%
  \NCSUStatementBegin{Lemma}{#1}{NCSUSoftBlueTitle}{NCSUBlueBorder}%
  \itshape
}{%
  \NCSUStatementEnd
}

\DeclareDocumentEnvironment{proposition}{O{}}{%
  \NCSUStatementBegin{Proposition}{#1}{NCSUSoftBlueTitle}{NCSUBlueBorder}%
  \itshape
}{%
  \NCSUStatementEnd
}

\DeclareDocumentEnvironment{corollary}{O{}}{%
  \NCSUStatementBegin{Corollary}{#1}{NCSUSoftBlueTitle}{NCSUBlueBorder}%
  \itshape
}{%
  \NCSUStatementEnd
}

\DeclareDocumentEnvironment{example}{O{}}{%
  \NCSUStatementBegin{Example}{#1}{NCSUBoxGrayTitle}{black!55}%
}{%
  \NCSUStatementEnd
}

\DeclareDocumentEnvironment{remark}{O{}}{%
  \NCSUStatementBegin{Remark}{#1}{NCSUBoxGrayTitle}{black!45}%
}{%
  \NCSUStatementEnd
}